\section{Introduction}

Avant de présenter les différentes fonctionnalités et interfaces de notre application. Il paraît nécessaire de noter que nous avons commencé le développement, dans un premier temps, par construire le BackEnd avec les commandes CLI de l'application, qui constitue l'âme du projet. Ensuite, nous nous sommes focalisés sur le FrontEnd afin d'essayer de présenter les résultats de tout ce qui se passe derrière.

Dans cette partie, nous nous intéresserons aux captures d'écran de la plateforme, accompagnées par des explications démontrant l'utilité de chacune des interfaces.

\section{Présentation des interfaces graphiques}

\subsection{Synchronisation des factures}

Cette page permet une visualisation claire l'interface de synchronisation de toutes les types existe dans AdminIs et ici principalement pour les Factures. Il affiche le nom de la synchronisation spécifier et la date de dernière synchronisation.

\subsection{Synchronisation des Balances ISOC}

Cette page permet une visualisation claire l'interface de synchronisation de toutes les types existe dans AdminIs et ici principalement pour les Balances ISOC. Il affiche le nom de la synchronisation spécifier et la date de dernière synchronisation.

\subsection{Synchronisation des Balances IPM}

Cette page permet une visualisation claire l'interface de synchronisation de toutes les types existe dans AdminIs et ici principalement pour les Balances IPM. Il affiche le nom de la synchronisation spécifier et la date de dernière synchronisation.

De plus, une option de synchroniser toutes les synchronisations de toutes les types prêts pour la Delta Synchro.

\subsection{Synchronisation des Déclarations TVA}

Cette page permet une visualisation claire l'interface de synchronisation de toutes les types existe dans AdminIs et ici principalement pour les RIC DECLARATION TVA. Il affiche le nom de la synchronisation spécifier et la date de dernière synchronisation.

\subsection{Synchronisation des Collaborateurs}

Cette page permet une visualisation claire l'interface de synchronisation de toutes les types existe dans AdminIs et ici principalement pour les Collaborateurs. Il affiche le nom de la synchronisation spécifier et la date de dernière synchronisation.

\subsection{Synchronisation des Clients}

Cette page permet une visualisation claire l'interface de synchronisation de toutes les types existe dans AdminIs et ici principalement pour les Clients. Il affiche le nom de la synchronisation spécifier et la date de dernière synchronisation.

\subsection{Synchronisation des Time Sheets}

Cette page permet une visualisation claire l'interface de synchronisation de toutes les types existe dans AdminIs et ici principalement pour les Time Sheet (Time Registration). Il affiche le nom de la synchronisation spécifier et la date de dernière synchronisation.

\section{Remarques}

Dans toutes les interfaces nous avons voir que on a l'option de 'Lancer Delta Globale'. Cette option a pour objectif de lancer toutes les synchronisations de toutes les types existe dans le Delta Synchro.

\section{Conclusion}

Ce chapitre est consacré à la description de la mise en œuvre de nouveau système de synchronisation. En présentant les interfaces illustrant le travail fait sur l'application Power Team.
